% generator_I38.tex -- Prop I.38 (Theor.): triangles on equal bases between same parallels are equal.
\documentclass[12pt]{article}\usepackage{geometry}\geometry{a5paper,margin=1.5cm}
\usepackage[lining]{ebgaramond}\usepackage{amsmath}\usepackage{amssymb}\usepackage{byrne}
\def\mpPre{u := 1cm; textLabels := false;}
\begin{document}
\defineNewPicture[1/5]{%
  pair A,B,C,D,E,F,G,H,I,J,d[];numeric h;%
  h:=5/2u;d1:=(3/2u,0);d2:=(3/4u,-h);d3:=(7/3u,h);d4:=(-1/4u,-h);%
  A:=(0,0);B:=A shifted d1;C:=A shifted d2;D:=C shifted d1;%
  E:=C shifted d3;F:=D shifted d3;G:=E shifted d4;H:=F shifted d4;%
  I:=(xpart(A),ypart(C));J:=(xpart(F),ypart(C));%
  draw byPolygon(A,B,C)(byyellow);%
  draw byPolygon(B,C,D)(byred);%
  draw byPolygon(E,F,H)(byblack);%
  draw byPolygon(E,G,H)(byblue);%
  draw byLine(B,D,byblue,SOLID_LINE,REGULAR_WIDTH);%
  draw byLine(E,G,byred,SOLID_LINE,REGULAR_WIDTH);%
  byLineDefine(A,C,byblue,DASHED_LINE,REGULAR_WIDTH);%
  byLineDefine(F,H,byred,DASHED_LINE,REGULAR_WIDTH);%
  byLineDefine(A,F,byblack,DASHED_LINE,REGULAR_WIDTH);%
  draw byNamedLineSeq(0)(AC,AF,FH);%
  draw byLine(I,J,byblack,DASHED_LINE,REGULAR_WIDTH);%
  draw byLabelsOnPolygon(A,B,E,F,noPoint)(ALL_LABELS,0);%
  draw byLabelsOnPolygon(H,G,D,C,noPoint)(ALL_LABELS,0);%
}
\noindent\begin{minipage}[t]{0.45\linewidth}\centering\vspace{0pt}%
\drawCurrentPicture\end{minipage}\hfill
\begin{minipage}[t]{0.52\linewidth}\vspace{0pt}%
\textbf{Book~I.\enspace Prop.\ XXXVIII. Theor.}
\medskip\noindent\itshape
Triangles (\drawPolygon[bottom]{BCD} and \drawPolygon[bottom]{EGH})
on equal bases and between the same parallels are equal.
\bigskip\noindent\upshape\begin{center}
$\left.\begin{aligned}
\mbox{Draw }\drawUnitLine{AC}&\parallel\drawUnitLine{BD}\\
\mbox{and }\drawUnitLine{FH}&\parallel\drawUnitLine{EG}
\end{aligned}\right\}$~(pr.~I.31)\\[4pt]
$\drawPolygon{ABC,BCD}=\drawPolygon{EFH,EGH}$~(pr.~I.36);\\[4pt]
but $\drawPolygon{ABC,BCD}=\mbox{twice }\drawPolygon{BCD}$~(pr.~I.34),\\
and $\drawPolygon{EFH,EGH}=\mbox{twice }\drawPolygon{EGH}$~(pr.~I.34),\\[4pt]
$\therefore \drawPolygon{BCD}=\drawPolygon{EGH}$~(ax.~VII).
\end{center}
\begin{flushright}Q.\ E.\ D.\end{flushright}
\end{minipage}
\end{document}
